\begin{pastexam}[label={px:ncku-113-ee-06}]{成大電機所 113 學年度 是非題第 6 題}
（是非；每題 3 分，答「錯」要指出原因，答錯倒扣 1 分）

Kruskal algorithm always keep a tree from the beginning to the end while it tries to construct a
minimum spanning tree.
\end{pastexam}

\begin{pastsoln}
這份考卷和成大電通所 113 學年度的「資料結構」是同一份（電機系、電通所、智慧資安學程共用），完整詳解見歷屆試題~\ref{px:ncku-113-cce-06}。

錯。 Kruskal 演算法過程中維持的是一個森林，不一定是一棵樹；一直保持一棵樹的是 Prim 演算法。
\end{pastsoln}
